\section{한계와 결과 해석 시 주의점}
\label{sec:threats}

\begin{revision}

\subsection{정답과 독립성}
기존 40쌍은 기대 결과를 확인해 남긴 규칙 회귀시험이다. 새 144건은 최초 결과에 따라 선별하지 않았지만 원문ㆍ프레임ㆍ기대 라벨과 비교기를 같은 설계자가 만들었다. 48개 실패를 사용한 조건별 모델은 같은 사례로 재시험했다. 개발 미사용 독립 자료, 외부 사람 의미 정답과 자연 오류 평가를 확보하지 못했다. 따라서 완전 검출, 독립 검증 완료나 자연 precision/recall/F1을 주장하지 않는다.

\subsection{자동 추출과 의미 표현}
조건별 모델은 자동 NLP가 아니다. 원문의 증거를 조건 ID로 올바르게 연결하는 단계가 작성 입력에 의존한다. 기존 파서는 제한된 영어 문구ㆍ명시적 연결사를 처리하고 관측 증거를 미상으로 둔다. 부정, 암시, 지시어, 복합행동, 복수 의무와 다른 언어에 대한 일반적 의미 추출은 검증하지 않았다. 단일 의무 투영 손실과 대기 중 누적 미상이 최초 결과를 낮췄다. 보완 모델은 모든 P1--P4 고장이나 임의의 다중 선행조건을 지원한 모델이 아니다.

\subsection{증거의 유효성과 모델 범위}
증거 보존은 같은 대상ㆍ조건과 명시한 구간이 유지된다는 가정을 쓴다. 객체 추적, 관측 노후화ㆍ만료, 해제 후 재등장과 시간에 따라 변하는 유효성은 포함하지 않았다. 복수 동시 문제의 사유 플래그를 보존하지만 대표 최초 의무는 정의한 우선순위로 하나를 선택한다. 이 출력이 모든 의무의 완전한 인과 설명은 아니다. 중심 모델은 이산 사건 순서이며 정량 시간 제한ㆍ브레이크/조향 명령ㆍ차량 동역학ㆍ폐루프 상호작용이 없다. 사건 끝에서 미래 해제를 못 봤다는 사실만으로 위반을 선언하지 않는다.

\subsection{자연 자료와 재현성}
27장면은 행동 표현으로 고른 비무작위 표본이다. 텍스트 불일치 23/27과 합의 사유 메모 0/27이 남아 있고, 검토자 전문성ㆍ교육ㆍ중재 및 논의 절차를 충분히 복원하지 못했다. 사람 재평가를 실시하지 않았다. 원 생성 실행별 모델 revision, prompt/config와 품질 점수 상류 임계값이 미복원이라 완전 재생성을 주장할 수 없다. 현재 입력ㆍ코드ㆍXMLㆍ질의ㆍ로그ㆍ해시는 보존했지만 restricted 자료의 외부 공개나 독립 사용 이력을 인증한 것은 아니다.

\subsection{보조 반응 분석과 일반화}
속도 변화는 실제 제어 명령의 대용이며 객체 해제의 참ㆍ거짓을 제공하지 않는다. $W,\delta,\rho,D$ 및 이미 정지 규칙은 탐색적 설정이고 안전 법규에서 유도하지 않았다. 27설정은 가능한 모든 조건을 포함하지 않는다. 궤적 불일치 7건을 자연 CoC 모순 정답으로 바꿀 수 없다. 다른 자료ㆍ언어ㆍ생성 모델로의 일반화, 실제 차량 안전과 학습 성능 개선은 미검증이다.
\end{revision}
